% !TeX root = ../../../thesis.tex

\subsubsection{\travisci}
\label{sssec:travisci}
The final CI platform that will be discussed is \travisci{}. This CI system was launched in 2011 and is only compatible with \github{} repositories. \travisci{} build tasks are configured similarly to GitLab CI but lack support for self-hosted runners. In addition to commit-triggered builds, we can also schedule daily, weekly or monthly builds using \emph{cronjobs}. Similar to \githubactions{}, open-source projects can use this service at zero cost, and a paid plan exists for private repositories \cite{travisanalysis}. It is not possible to create custom plugins, but \travisci{} already features built-in support for a variety of programming languages. In 2020, almost one million projects are using \travisci{} \cite{travis2020}.

\begin{figure}[htbp!]
	\centering
	\includegraphics[width=0.40\textwidth]{assets/images/travis-ci.pdf}
	\caption{Logo of \travisci{} (\url{https://travis-ci.com/}).}
	\label{fig:travis-ci}
\end{figure}
